\section{Harmonic extensions and finite balance}\label{sec:extensions}

We first identify the algebra that enters the finite model. The use of
the complete algebra, rather than one scalar per stable type, determines
both the number of small eigenvalues and the correct stability condition.

\begin{lemma}\label{lem:algebra}
For the extension \eqref{eq:extension}, harmonic representatives of its
fixed cohomology class give a presentation \eqref{eq:harmonic} for every
small $s$. In this presentation
\begin{equation}\label{eq:centralizer}
H^0(X,\End E)=
\left\{(X_v)\in\bigoplus_vM_{m_v}(\C):
X_iZ^0_{ij,a}=Z^0_{ij,a}X_j\text{ for all }i,j,a\right\}.
\end{equation}
The two algebras are identified by parallel block endomorphisms.
Every holomorphic isomorphism between two harmonic presentations is a
parallel block gauge in $\mathcal G$. In particular simplicity of $E$
is equivalent to the common centralizer being scalar.
\end{lemma}

\begin{proof}
Stable nonisomorphic bundles of a common slope have no nonzero maps
between them, and each has only scalar endomorphisms; see
\cite[Chapter~V, \S7]{kobayashi1987}. Hence
$\Hom(F_v,F_w)=0$ for $v\ne w$ and
$\End(F_v)=\C\id$. These are statements about the fixed holomorphic
bundles and do not change with $s$.

A smooth splitting of \eqref{eq:extension} gives a closed off-diagonal
form $\eta\in\Omega^{0,1}(\Hom(U,S))$. Its harmonic part is
$N_s(Z^0)$, and $\eta-N_s(Z^0)=\dbar C_s$ for a smooth section
$C_s$ of $\Hom(U,S)$. An upper-triangular unipotent change of splitting
removes this exact part. Products of two source-to-target blocks vanish,
so there is no nonlinear term in this change. The same observation shows
that $\dbar_{s,Z}^2=0$ for every coefficient array $Z$. Taking $C_s$
orthogonal to the kernel gives a smooth choice in $s$. This gauge has
determinant one.

Apply $\Hom(F_i,-)$ to \eqref{eq:extension}. Since
$\Hom(F_i,U)=0$, one obtains
$\Hom(F_i,E)=\Hom(F_i,S)=\C^{m_i}$. The evaluation image
\[
\operatorname{im}\left(\bigoplus_{i\in I}
F_i\otimes\Hom(F_i,E)\longrightarrow E\right)=S
\]
is the displayed locally free subbundle, fiber by fiber. Every
endomorphism preserves it. In a harmonic splitting an endomorphism
therefore has the form $\left(\begin{smallmatrix}A&C\\0&B\end{smallmatrix}\right)$.
Its diagonal equations make $A$ and $B$ parallel multiplicity blocks.
Its remaining equation is
\begin{equation}\label{eq:exact-harmonic}
\dbar C+N_s(Z^0)B-AN_s(Z^0)=0.
\end{equation}
The last two terms are harmonic, since the multiplicity blocks are
parallel. Exact and harmonic forms are orthogonal, so each term in this
decomposition vanishes. Now $C$ is holomorphic and
$\Hom(U,S)=0$, whence $C=0$. Independence of the cohomology bases gives
\eqref{eq:centralizer}. Conversely those equations plainly define a
holomorphic endomorphism. The same proof for an isomorphism between two
presentations gives the asserted gauge statement.
\end{proof}

The separation of stable types across the layers is natural in this
simple two-step setting. Indeed, for any extension of two polystable
layers of equal slope the map
\[
\Hom(U,S)\longrightarrow\End(E),\qquad f\longmapsto\iota f\pi
\]
is injective and has square-zero image. Simplicity forces this space to
vanish. If a stable type occurred in both layers, it would give a nonzero
map from $U$ to $S$. This observation is specific to two polystable
layers and is not a claim about longer filtrations.

\subsection{The real domain and its dual}
The Hermitian logarithms used for metric corrections form
\begin{equation}\label{eq:realspaces}
\hh=\left\{Y_v=Y_v^*: \sum_vr_v\tr Y_v=0\right\},\qquad
\mathcal W=\left\{T_v=T_v^*: \sum_v\tr T_v=0\right\}.
\end{equation}
They are paired by $\langle Y,T\rangle=\sum_v\tr(Y_vT_v)$.
The positive definite real inner product on $\hh$ is
$\langle X,Y\rangle_r=\sum_vr_v\tr(X_vY_v)$; division of the
$v$-th block of $T$ by $r_v$ identifies $\mathcal W$ with $\hh$.
Both $\mu_s(Z)$ and $b=(b_vI_{m_v})$ lie in $\mathcal W$.
Multiplicities $m_v$ and factor ranks $r_v$ have different roles.

For example, at a fixed parameter orthonormalize the harmonic arrow
basis in the volume-normalized inner product. Then
\begin{equation}\label{eq:momentcomponents}
\mu_i(Z)=\sum_{j,a}Z_{ij,a}Z_{ij,a}^*,\qquad
\mu_j(Z)=-\sum_{i,a}Z_{ij,a}^*Z_{ij,a}.
\end{equation}
The rank weights have not disappeared: they occur in $b_v$ and in the
inner product on the endomorphism algebra. In the original cohomology
coordinates the same equations use the corresponding positive Gram
matrices. Changing arrow basis commutes with $\mathcal G$ and preserves
all subrepresentations.

\begin{proposition}\label{prop:balance}
If the fixed representation $Z^0$ is $b$-stable, its orbit contains a
balanced representative. Under the scalar-centralizer hypothesis this
representative is unique up to $K$. For every such representative
\begin{equation}\label{eq:amplitude}
q_s(Z)=\sum_{i\in I}m_ib_i=\sigma.
\end{equation}
The stable real weights for this fixed representation form an open
polyhedral cone in $\{b:\sum_vm_vb_v=0\}$, possibly with empty
intersection with the accessible parameter image.
\end{proposition}

\begin{proof}
The real-weight finite Kempf--Ness theorem
\cite[\S5.1, Theorem~5.1]{hkkp2023} states that a representation admits
a balanced metric precisely when it is polystable. Its weight convention
is $\theta(W)\le0$; setting $\theta_v=-b_v$ gives
\eqref{eq:finite-stability} and \eqref{eq:momentcomponents}.
The positive masses in the group metric are $r_v$. Equivalently the
central endomorphism parameter of that convention is $-b_v/r_v$.
No second rank factor enters \eqref{eq:finite-stability}.
The theorem allows real weights through its logarithmic determinant
potential; no rational approximation is required. The underlying
integral-character statement is King's
\cite[Propositions~3.1--3.2 and Theorem~6.1]{king1994}.

For clarity, uniqueness here can be checked directly. If two balanced
points differ by $g\in\mathcal G$, polar decomposition and a common
scalar give $g=ke^{Y/2}$, with $k\in K$ and $Y\in\hh$.
Diagonalize $Y$ blockwise and denote its eigenvalues by $y_p$.
Writing $w_{pq}$ for the total squared arrow entries from source
eigenvector $q$ to target eigenvector $p$, the function
\begin{equation}\label{eq:convexpotential}
f(t)=q_s(e^{tY/2}\cdot Z)-t\langle Y,b\rangle
\end{equation}
satisfies
\[
f''(t)=\sum_{p,q}(y_p-y_q)^2e^{t(y_p-y_q)}w_{pq}.
\]
For $Y\ne0$ in $\hh$ this is positive: otherwise $Y$ centralizes
$Z$ and must be scalar, contrary to its weighted trace zero.
Balance at both endpoints would give $f'(0)=f'(1)=0$.
Thus $Y=0$. Tracing all target equations in
\eqref{eq:momentcomponents} proves \eqref{eq:amplitude}.
Finally only finitely many dimension vectors $(\dim W_v)$ are
possible, and each gives a linear strict inequality. These define the
claimed cone. The actual polarization image need not fill that cone.
\end{proof}

At balance the sign can also be seen without a stability convention.
If $\Pi_v$ projects onto a subrepresentation, then for each arrow
\[
\tr(\Pi_iZZ^*)-\tr(\Pi_jZ^*Z)
=\|\Pi_iZ(1-\Pi_j)\|_F^2-\|(1-\Pi_i)Z\Pi_j\|_F^2.
\]
The second term is zero. Equality after summation would make both
the subrepresentation and its orthogonal complement invariant; its
projection would then be a nonscalar centralizer element. Thus balance
with scalar centralizer implies precisely the strict positive sign.

A common complex scalar can normalize the bundle determinant of any
balancing gauge to one, since it acts trivially on the arrows and its
bundle determinant is its $R_E$-th power. Choices of scalar phase do
not affect the metric. Together with the unipotent identification in
Lemma~\ref{lem:algebra}, this defines the determinant convention in
Theorem~\ref{thm:main}.
